\chapter{Modal Logic}
\label{ch:modal}

Modal logic adds operators — canonically $\Box$ (necessity) and
$\Diamond$ (possibility) — that make a statement's truth depend on
which \emph{state} it is evaluated at, connected to other states by
an accessibility relation. This chapter introduces the accessibility
picture directly through a structure already familiar from outside
logic proper: a fact moving from merely recorded, to admitted, to
committed, to consequential, is a chain of states each reachable
from the last, and modal vocabulary describes exactly that kind of
movement.

\section{States and Accessibility}

Formally, $\Box p$ at a state $w$ means ``$p$ holds at every state
accessible from $w$,'' and $\Diamond p$ means ``$p$ holds at some
state accessible from $w$.'' The accessibility relation itself is
nothing exotic — an ordinary relation on a type of states, exactly
the kind of object Chapter~\ref{ch:admissibility} builds
\texttt{Extends} out of.

\begin{experiment}{A three-world Kripke frame, evaluated by hand}
Three worlds, a chain of accessibility $w_0 \to w_1 \to w_2$, and a
proposition $p$ true only at $w_1$.
\begin{lean}
inductive World where
  | w0 | w1 | w2
  deriving DecidableEq, Repr

def accessible : World → World → Bool
  | .w0, .w1 => true
  | .w1, .w2 => true
  | _, _ => false

def val : World → Bool
  | .w0 => false
  | .w1 => true
  | .w2 => false

def worlds : List World := [World.w0, World.w1, World.w2]

def diamondP (w : World) : Bool :=
  worlds.any (fun v => accessible w v && val v)

def boxP (w : World) : Bool :=
  worlds.all (fun v => !accessible w v || val v)
\end{lean}
By hand, before automating: at $w_0$, the only accessible world is
$w_1$, where $p$ holds — so both $\Diamond p$ and $\Box p$ hold at
$w_0$ (there is nothing accessible for $\Box$ to fail on). At $w_2$,
nothing is accessible at all, so $\Diamond p$ fails outright (there
is no witness) while $\Box p$ holds \emph{vacuously} — a dead-end
world satisfies every necessity claim for free.
\begin{lean}
example : diamondP World.w0 = true := by decide
example : boxP World.w0 = true := by decide
example : diamondP World.w2 = false := by decide
example : boxP World.w2 = true := by decide
\end{lean}
\end{experiment}

\section{A Worked Structure: Record, Admit, Commit}

\begin{experiment}{The forward chain}
A fact can be recorded without being admitted, admitted without
being committed, and committed without being consequential — but
never the reverse. State this as a chain of implications between
four propositional fields of one structure.
\begin{lean}
structure State where
  recorded : Prop
  admitted : Prop
  committed : Prop
  consequential : Prop

def valid (s : State) : Prop :=
  (s.admitted → s.recorded) ∧
  (s.committed → s.admitted) ∧
  (s.consequential → s.committed)

theorem consequence_requires_record
    (s : State) (h : valid s) (hc : s.consequential) :
    s.recorded :=
  h.1 (h.2.1 (h.2.2 hc))
\end{lean}
Read modally: if ``consequential'' is accessible from
``committed,'' ``committed'' from ``admitted,'' and ``admitted''
from ``recorded,'' then $\Box$ at the recorded state reaches every
later state in the chain — but nothing forces the reverse
accessibility.
\end{experiment}

\begin{countermodel}{None of the converses hold}
Recording does not imply admission; admission does not imply
commitment.
\begin{lean}
def recordedOnly : State where
  recorded := True
  admitted := False
  committed := False
  consequential := False

theorem record_does_not_imply_admit :
    ¬ ∀ s : State, valid s → s.recorded → s.admitted := by
  intro h
  exact h recordedOnly (by simp [valid, recordedOnly]) (by simp [recordedOnly])
\end{lean}
\texttt{recordedOnly} is a single state satisfying \texttt{valid}
in which every later stage of the chain is false — exactly the kind
of witness a countermodel to a converse needs, and exactly the
reason modal accessibility is not assumed symmetric by default.
\end{countermodel}

\section{Necessity, Possibility, and Frame Conditions}

Standard modal axioms correspond exactly to conditions on the
accessibility relation itself, a fact usually called
\emph{correspondence theory}: axiom T ($\Box p \to p$) corresponds
to \emph{reflexivity} (every world sees itself, so what necessarily
holds already holds here); axiom 4 ($\Box p \to \Box\Box p$)
corresponds to \emph{transitivity} (if every step-1-accessible world
satisfies $p$, so does every step-2-accessible world, since step-2
accessibility factors through step-1); axiom B ($p \to \Box\Diamond p$)
corresponds to \emph{symmetry} (if $w$ sees $v$, $v$ sees $w$ back,
so whatever holds at $w$ is reachable-back-to from anywhere $w$
reaches). The three-world frame above is not transitive as drawn —
$w_0$ reaches $w_2$ only through $w_1$, not directly — which is
exactly why axiom 4 is not being claimed for it.

The record/admit/commit chain, by contrast, \emph{is} transitive in
exactly the relevant sense, and the composition already proved in
\texttt{consequence\_requires\_record} is the modal axiom~4 pattern
made concrete: reaching \texttt{recorded} from \texttt{consequential}
factors through the intermediate states, the same way
\texttt{extends\_trans} in Chapter~\ref{ch:admissibility} composes
two one-step extensions into a two-step one.

\begin{experiment}{Two-step accessibility, factored}
\begin{lean}
theorem committed_reaches_recorded (s : State) (h : valid s) :
    s.committed → s.recorded :=
  fun hc => h.1 (h.2.1 hc)
\end{lean}
Compare the shape directly to \texttt{extends\_trans}'s
\texttt{fun x hx => h23 (h12 hx)}: both compose two one-step
implications into one two-step implication by applying them in
sequence. \texttt{h.2.1} is the one-step link
\texttt{committed → admitted}; \texttt{h.1} is the next one-step
link \texttt{admitted → recorded}; composing them is the entire
proof, and it is the same composition — accessibility chaining
through an intermediate state — in both settings.
\end{experiment}

\section{Exercises}

\begin{exercise}
Show that admission does not imply commitment, following the pattern
of \texttt{record\_does\_not\_imply\_admit} above.
\end{exercise}
